\section{The Cosmos: Expansion and the Dark Sector}

Step all the way back from particles and look at the whole universe. A universe that is a growing crystal should have things to say about cosmology: why it expands, where its arrow comes from, and the two great mysteries of modern astronomy, dark energy and dark matter. This chapter sketches what the crystal says.

\subsection{Expansion comes for free}

The universe is expanding. Galaxies drift apart, space itself stretches. In ordinary cosmology this is an observed fact fed into the equations. In the crystal it is almost a tautology. The crystal grows, adding cells at its frontier forever. A growing web of relations, seen from inside, just is an expanding space. You do not have to arrange for expansion. A universe that grows expands by definition.

The simulation shows the expected shape: a brief, dramatic early burst of fast growth, what cosmologists call inflation, followed by a settling into steadier expansion. And the arrow of time, the one-way direction we built in the bootstrap, is the same one-way growth, so the universe's expansion and its arrow of time are the very same fact, not two coincidences.

\subsection{Dark energy: the everpresent constant}

The deepest mystery in cosmology is dark energy, a faint pressure pushing the expansion to speed up. Its value is bizarre: not zero, but staggeringly tiny, smaller than naive physics predicts by a number with over a hundred zeros. Why so small but not zero? This is often called the worst prediction in physics.

The crystal offers a clean idea. The strength of this push is not a fixed constant set once at the start. It is the natural jitter in the size of spacetime, and that jitter shrinks as the universe grows, falling like one over the square root of the total volume. Plug in the actual size of our universe, which is vast, and the jitter comes out tiny, landing remarkably close to the observed dark energy, to the right rough size, with nothing tuned. The smallness is not a fluke to be explained away. It is what you get because the universe is huge and the push is a shrinking fluctuation. This is one of the theory's most striking contacts with real numbers.

\subsection{Dark matter: a candidate}

The other mystery is dark matter: galaxies spin as though far more mass is present than we can see, holding their outer stars from flying off. The crystal has a candidate mechanism, a small long-range adjustment to how gravity behaves on galactic scales, arising from the structure of the web, that flattens the spin of galaxies without inventing a new invisible particle. This one is more tentative than the dark energy result, a plausible mechanism rather than a sharp prediction, and the walkthrough flags it as such. But it shows the crystal has resources to address the problem from its own structure.

\subsection{The shape of the cosmic story}

Put together, the crystal tells a coherent cosmic story. A seed unfolds by the growth rule. There is a fierce early burst, then a graceful expansion, with an arrow of time that is simply the growth. The acceleration we see today is the shrinking jitter of a now-enormous universe, which is why it is so faint. And the missing galactic mass may be a structural quirk of gravity on the web rather than new stuff. Some of this is solid (expansion, the arrow, the dark-energy size), some is a promising first step (dark matter, the fine details of structure forming). The next chapters narrow to the smallest scales and to what could be tested.

\textbf{Takeaway:} the crystal expands simply because it grows, with an early inflationary burst, and its expansion and arrow of time are one fact. Dark energy comes out tiny but nonzero as the shrinking jitter of a vast spacetime, matching the observed value in size, and dark matter has a candidate structural mechanism that needs no new particle.
